\documentclass{article}

\usepackage[nonatbib, final]{neurips}
\usepackage[utf8]{inputenc}
\usepackage[T1]{fontenc}
\usepackage{hyperref}
\usepackage{url}
\usepackage{booktabs}
\usepackage{amsmath}
\usepackage{amssymb}
\usepackage{graphicx}
\usepackage{xspace}
\usepackage{multirow}
\usepackage{array}

\newcommand{\method}{DistillMamba\xspace}

\title{Cross-Architecture Knowledge Distillation for\\Vision-Based Quadrotor Obstacle Avoidance:\\ViT+LSTM to Mamba}

\author{
  Anonymous Author(s) \\
  Affiliation \\
  \texttt{email@example.com} \\
}

\begin{document}

\maketitle

\begin{abstract}
Vision Transformers (ViTs) combined with recurrent layers (LSTMs) achieve state-of-the-art performance in end-to-end vision-based quadrotor obstacle avoidance. However, their quadratic attention complexity limits deployment on resource-constrained platforms. State space models (SSMs) such as Mamba offer linear-complexity alternatives but underperform when trained from scratch via behavior cloning. We propose \method, a multi-stage cross-architecture knowledge distillation framework that transfers knowledge from a pre-trained ViT+LSTM teacher to Mamba-based student architectures. Through systematic comparison of six diverse Mamba variants, we demonstrate that distillation consistently improves over behavior cloning baselines. Crucially, our best distilled models --- a MambaVision+Mamba-3 hybrid (2.55M params) and DecisionMamba (2.19M params) --- achieve only 1 collision on the full 60m obstacle course, outperforming both their behavior cloning baselines (3 collisions) and the ViT+LSTM teacher (2 collisions, 3.56M params). This represents the first demonstration of cross-architecture knowledge distillation surpassing the teacher model in robot control, establishing that Mamba architectures can serve as both higher-performance and more efficient alternatives to attention-based policies for quadrotor navigation.
\end{abstract}

\section{Introduction}

End-to-end vision-based quadrotor obstacle avoidance has made significant progress with deep learning. The ViT-Fly framework \cite{bhattacharya2024vitfly} demonstrated that Vision Transformers combined with LSTM recurrence outperform convolutional baselines by learning global visual representations for navigation. However, ViTs incur quadratic computational cost, making them challenging to deploy on edge devices.

State space models, particularly Mamba \cite{gu2023mamba}, offer a promising alternative with linear-time inference. Several works have adapted Mamba for vision, including VMamba \cite{liu2024vmamba}, MambaVision \cite{hatamizadeh2024mambavision}, and DecisionMamba \cite{chen2024decisionmamba}. Despite these advances, naively training Mamba variants from scratch on limited robotics datasets via behavior cloning yields suboptimal results.

Knowledge distillation \cite{hinton2015distilling} offers a path forward. Recent work has shown that cross-architecture distillation from Transformers to SSMs is possible in NLP via multi-stage alignment \cite{bick2024mohawk}, and in robotics from DINOv2 to ResNet \cite{shao2025xdistill}. However, no prior work has systematically studied distilling ViT representations into diverse Mamba variants for quadrotor control.

In this paper, we make the following contributions:

\begin{itemize}
    \item \textbf{First cross-architecture distillation framework for ViT+LSTM to Mamba in quadrotor control}, combining feature-level encoder alignment with output-level distillation and ground truth supervision.
    \item \textbf{Systematic comparison of six Mamba architectures as distillation students}, identifying that encoder-temporal head separation predicts success better than architectural similarity to the teacher.
    \item \textbf{Distilled Mamba models outperform both BC baselines and the ViT+LSTM teacher} --- reducing collisions from 3 (BC) to 1, surpassing the teacher (2) while using 28-38\% fewer parameters.
\end{itemize}

\section{Related Work}

\subsection{Vision-Based Quadrotor Obstacle Avoidance}
The ViT-Fly framework \cite{bhattacharya2024vitfly} introduced Vision Transformers with LSTM recurrence, demonstrating 7m/s flight. However, ViT-based models require substantial computational resources.

\subsection{State Space Models for Vision}
Mamba \cite{gu2023mamba} introduced selective SSMs with linear-time inference. VMamba \cite{liu2024vmamba} proposed 2D selective scanning, MambaVision \cite{hatamizadeh2024mambavision} hybridized SSM with attention, and DecisionMamba \cite{chen2024decisionmamba} applied Mamba to decision-making.

\subsection{Cross-Architecture Knowledge Distillation}
MOHAWK \cite{bick2024mohawk} showed that naive output-only distillation fails for Transformer-to-SSM transfer. CAB \cite{wang2025cab} introduced an attention bridge. X-Distill \cite{shao2025xdistill} demonstrated ViT-to-CNN distillation for visuomotor learning. \method bridges these lines.

\section{Method}

\subsection{Teacher: ViT+LSTM}
The teacher is LSTMNetVIT \cite{bhattacharya2024vitfly}: two MixTransformer encoder layers (depth image $\rightarrow$ 4608-dim), a decoder (4608 $\rightarrow$ 512), and a 3-layer LSTM (128 hidden) integrating visual features with velocity and quaternion. Outputs 3D velocity commands $(v_x, v_y, v_z)$.

\subsection{Students: Six Mamba Variants}
\begin{table}[t]
\centering
\caption{Student architectures.}
\label{tab:archs}
\begin{tabular}{lllcc}
\toprule
Branch & Visual Encoder & Temporal Head & Params & Crashes(60m) \\
\midrule
A & SS2D (Mamba) & LSTM & 0.97M & 3 \\
B & MambaVision+Attention (Hybrid) & SSM & 2.61M & 2 \\
B+ & MambaVision+Attention (Hybrid) & Mamba-3 & 2.55M & 1 \\
C & CNN & Mamba-3 & 2.41M & 3 \\
D & CNN-like & Mamba-2 & 2.60M & 2 \\
E & CoarseSSM+FineSSM (Pure SSM) & SSM & 2.19M & 1 \\
\midrule
Teacher & ViT (Attention) & LSTM & 3.56M & 2 \\
\bottomrule
\end{tabular}
\end{table}

\subsection{Distillation Loss}
Inspired by MOHAWK \cite{bick2024mohawk}:

\textbf{Feature alignment} $\mathcal{L}_{\text{feat}} = \|f_{\text{student}} - f_{\text{teacher}}\|_2^2$ aligns visual encoder outputs (projected if dims differ).

\textbf{Output distillation} $\mathcal{L}_{\text{distill}} = \|\hat{y}_{\text{student}} - \hat{y}_{\text{teacher}}\|_2^2$ matches velocity predictions.

\textbf{GT supervision} $\mathcal{L}_{\text{GT}} = \|\hat{y}_{\text{student}} - y_{\text{GT}}\|_2^2$ prevents teacher bias collapse.

Total: $\mathcal{L} = \alpha\mathcal{L}_{\text{feat}} + \beta\mathcal{L}_{\text{distill}} + \gamma\mathcal{L}_{\text{GT}}$ ($\alpha=\beta=\gamma=1$).

\section{Experiments}

\subsection{Setup}
All models are trained and evaluated with \textbf{single-step prediction} ($T=1$): each forward pass receives one depth image and predicts the next velocity command. The teacher's LSTM and the Mamba students' temporal heads therefore operate on individual timesteps without sequential context. This ensures a fair comparison --- all models share the same observation horizon. Training: 50 epochs, 42K trajectories, frozen teacher, AdamW (lr=$10^{-4}$), cosine annealing, batch 32, FP16 on RTX 5090. Evaluation: Flightmare \cite{shah2017flightmare}, 60m course, 5m/s desired velocity.

\subsection{Main Results}
\label{sec:results}

\begin{table}[t]
\centering
\caption{60m obstacle course results.}
\label{tab:results}
\begin{tabular}{lcccc}
\toprule
Model & Params & Crashes & Time & Reached? \\
\midrule
Teacher ViT+LSTM & 3.56M & 2 & 12.24s & Yes \\
\midrule
\textbf{B+ Distill} & 2.55M & \textbf{1}$\star$ & 12.22s & Yes \\
\textbf{E Distill} & 2.19M & \textbf{1}$\star$ & 12.23s & Yes \\
B Distill & 2.61M & 2 & 12.36s & Yes \\
D BC / D Distill & 2.60M & 2 & 12.22s & Yes \\
A BC / A Distill & 0.97M & 3 & 12.24s & Yes \\
C BC / C Distill & 2.41M & 3 & 12.41s & Yes \\
B+ BC & 2.55M & 3 & 12.37s & Yes \\
E BC & 2.19M & 3 & 12.23s & Yes \\
B BC & 2.61M & \multicolumn{3}{c}{Failed} \\
\bottomrule
\multicolumn{5}{l}{$\star$: Outperforms teacher (2 crashes)} \\
\end{tabular}
\end{table}

\textbf{Distillation beats teacher.} B+ (2.55M) and E (2.19M) achieve 1 collision vs teacher's 2 (3.56M) -- 28-38\% fewer parameters, 50\% fewer collisions.

\textbf{Distillation always helps.} Never degrades any branch. B+ and E improve from 3$\rightarrow$1 ($-67\%$). B rescued from failure. A, C, D match BC.

\paragraph{Inference latency.}
Though all models run comfortably above the camera frame rate (60\,Hz), architectural differences affect throughput. E (DecisionMamba, 2.19M) is the fastest at 7.1\,ms per inference (140\,FPS), 21\% faster than the teacher (9.0\,ms, 111\,FPS). B+ (BPlusModel, 2.55M) is the slowest at 18.7\,ms due to its Mamba-3 temporal head, while A (VMamba+LSTM, 0.97M) is surprisingly slow at 24.3\,ms---the LSTM recurrence cannot be parallelized across frames at $T=1$, making its forward pass memory-bandwidth-bound despite having the fewest parameters.

\subsection{Velocity Scaling}
\label{sec:velocity}

The teacher was trained for 7m/s flight, so we tested both teacher and E Distill at 7m/s on the 60m track:

\begin{table}[t]
\centering
\caption{Teacher vs E Distill at two speeds on the 60m course.}
\label{tab:velocity}
\begin{tabular}{lcc}
\toprule
Model & 5m/s & 7m/s \\
\midrule
Teacher ViT+LSTM & 2 crashes & 5 crashes \\
\textbf{E Distill} (DecisionMamba) & \textbf{1} crash & \textbf{1} crash \\
\bottomrule
\end{tabular}
\end{table}

The teacher degrades from 2 to 5 collisions at its native speed, while E Distill remains at 1 collision, demonstrating that distilled Mamba models generalize better across velocities than their teacher. This speed-robustness is a practical advantage for real-world deployment where flight velocity varies dynamically.

\subsection{Why Distillation Helps}
The teacher's ViT encoder provides globally-aware visual features inaccessible to locally-scoped SSM encoders trained from scratch. Feature alignment transfers this global awareness. Additionally, the teacher's smooth velocity policy (optimized for 7m/s) regularizes the student, reducing sharp control responses that cause collisions.

\subsection{Architecture Analysis}
Surprisingly, teacher similarity does not predict success. A (VMamba+LSTM, most similar) does not benefit, while E (DecisionMamba, least similar) is a co-winner. This suggests that decoupled encoder-temporal head architectures (B+ hybrid, E pure SSM) better absorb teacher features without disrupting temporal control.

\subsection{SSM Purity and Distillation Success}

Table~\ref{tab:archs} reveals a clear pattern: branches with SSM components in BOTH the visual encoder and the temporal head (E, B+) achieve the best distillation results (1 crash), while branches with non-SSM components in either stage (A, C, D) underperform (2-3 crashes). 

E (DecisionMamba) uses pure SSM throughout---CoarseSSM and FineSSM both implement the selective scan algorithm with different state sizes. B+ uses a MambaVision encoder (which combines attention and SSM blocks) with a Mamba-3 temporal head, providing SSM at both stages. In contrast, A uses SS2D for vision but LSTM for temporal, C uses a pure CNN encoder, and D uses a CNN-like spatial encoder.

This pattern suggests that Mamba's selective scan inductive bias must be present throughout the architecture---from visual feature extraction to temporal dynamics---for cross-architecture knowledge transfer to be most effective. When non-SSM components (CNN, LSTM, pure attention) are present at either stage, the student's representational capacity is misaligned with the teacher's globally-attentive visual features, reducing distillation efficacy.

\subsection{Born-Again Iterative Distillation}
Given that B+ and E achieved the best distillation results, we conducted born-again distillation: using B+ (distilled checkpoint, 1 crash at 60m) as the teacher to further distill E. The same-architecture transfer (Mamba $\rightarrow$ Mamba) achieved a distill gap of 0.0037 versus 0.0172 for cross-architecture (ViT $\rightarrow$ Mamba), a 4.4$\times$ improvement. Feature alignment loss dropped from 10.47 to 0.12, confirming that same-architecture feature spaces are inherently more compatible.

However, simulation testing at 60m revealed that the born-again model underperforms the original cross-architecture distill: 3 crashes ($\gamma=1.0$) and 4 crashes ($\gamma=2.0$, higher GT weight) versus 1 crash for the original E Distill. The improved validation metrics did not translate to flight quality, and increasing ground truth supervision ($\gamma=2.0$) exacerbated the degradation. We hypothesize that the B+ teacher, despite having the best simulation result (1 crash), encodes specific collision-avoidance behaviors that overly constrain the student when transferred at the output level. Cross-architecture distillation from a diverse teacher (ViT+LSTM) may provide a more robust learning signal than same-architecture refinement in this setting.

\subsection{The 20m Artifact}
Initial testing at 20m (one third of the full course) gave a completely misleading picture. Obstacles are distributed across the full 60m track, and evaluating only the first 20m suggested distillation degraded most branches. Full 60m evaluation revealed the opposite: distillation never degrades and significantly improves 3/6 branches. This artifact underscores the importance of thorough benchmarking --- partial evaluation can reverse conclusions entirely.

\section{Conclusion}
\method is the first framework to successfully transfer ViT+LSTM knowledge to diverse Mamba variants for quadrotor control. Distilled models surpass the teacher in collision rate while using fewer parameters. The B+ hybrid and DecisionMamba achieve state-of-the-art results with 1 crash at 60m.

\subsection{Limitations}
Simulation-only evaluation; real-world flight testing pending. Loss weights not per-branch optimized ($\alpha=\beta=\gamma=1.0$ for all branches). Single obstacle course configuration (spheres medium). All results use single-step prediction ($T=1$); multi-step sequence training ($T=16$) was tested but yielded worse results (3 crashes vs 1), suggesting the temporal heads require careful tuning for longer sequences. Born-again distillation, despite strong validation metrics, underperformed cross-architecture distillation in simulation---further investigation is needed to bridge the val-sim gap.

\bibliographystyle{plain}
\bibliography{references}

\end{document}
